%% ============================================================
%% introduction.tex  —  Chapter 1
%% ============================================================
\chapter{Introduction}
\label{ch:introduction}

\section{Titan as an Astrobiological Target}
\label{sec:search_context}

The search for habitable environments beyond Earth has increasingly focused on icy ocean worlds in the outer solar system.  Among these, Saturn's moon \titan{} occupies a singular position.
It is the only moon with a dense atmosphere --- dominated by molecular nitrogen
at 1.467\,bar surface pressure and 93.65\,K \citep{Fulchignoni2005} --- and the
only known extraterrestrial world to host stable surface liquids
\citep{Stofan2007}, completing a full precipitation--evaporation--lake cycle
structurally analogous to Earth's hydrological cycle.  Its atmosphere
continuously produces complex organic tholins at
$\sim5\times10^{-14}$\,g\,cm$^{-2}$\,s$^{-1}$ \citep{Lavvas2011}, building
a global organic sediment layer over $\sim$4\,Gyr \citep{Lorenz2008}.
At depth, a global water--ammonia ocean at $\sim$55--80\,km is confirmed by the
tidal Love number $k_2 = 0.589\pm0.150$ \citep{Iess2012}.  This convergence of
cycling solvent, vast organic inventory, and subsurface water makes \titan{} a
natural laboratory for prebiotic chemistry operating at planetary scale
across geological time.

\section{The Physical Environment}
\label{sec:titan_environment}

The atmosphere ($\sim$94\% \ce{N2}, 5.65\% \ce{CH4}; \citealt{Niemann2005})
sustains a methane cycle with observed equatorial rainfall \citep{Turtle2011},
fluvial channel networks \citep{Miller2021channels}, and polar hydrocarbon
seas (Kraken, Ligeia, and Punga Mare) containing $\sim$70{,}000\,km$^3$ of
liquid methane--ethane \citep{Mastrogiuseppe2019,Malaska2025}.  The Lopes et al.\ \citeyearpar{Lopes2019} geomorphological map, derived from the full Cassini SAR dataset, classifies the surface into seven terrain units.  The surface is geologically diverse: equatorial dune fields ($\sim$17\% by SAR area), plains
($\sim$65\%), labyrinth terrain, impact craters, and mountains, as classified
by the Lopes et al.\ \citeyearpar{Lopes2019} global geomorphological map that
provides the terrain classification underpinning this thesis.

\section{Chemical Habitability Scenarios}
\label{sec:prior_theory}

Four habitability pathways have been proposed for \titan{}'s surface.
\textit{Acetylenotrophy} \citep{McKaySmith2005}: the reaction
$\ce{C2H2 + 3H2 -> 2CH4}$ ($\Delta G = -334$\,kJ\,mol$^{-1}$) provides
energy well above the minimum for terrestrial methanogens, even under updated
thermochemical estimates \citep{Yanez2024}.  Acetylene is supplied by UV
photolysis and detected at the surface by GCMS \citep{Niemann2005}.
\textit{Azotosome membranes} \citep{Stevenson2015}: acrylonitrile (detected by
ALMA; \citealt{Palmer2017}) spontaneously forms stable vesicle structures in
liquid methane, providing a route to cellular compartmentalisation without
aqueous chemistry.
\textit{Impact-melt prebiotic chemistry}: large impacts produce transient
water--ammonia--organic melt ponds persisting for $10^3$--$10^4$\,yr, during
which Strecker synthesis yields amino acids from surface aldehydes, ammonia,
and HCN \citep{Neish2010,Madan2026}.
\textit{Red-giant ocean}: when solar luminosity raises \titan{}'s surface
above the water--ammonia eutectic (176\,K; \citealt{Grasset2005}), a global
ocean will form over billions of years of accumulated organics for an
estimated 200--500\,Myr \citep{Lorenz1997}.  This thesis provides the first
unified quantitative characterisation of all four scenarios across
geological time.

\section{Scope and Research Questions}
\label{sec:scope}

Previous habitability assessments of \titan{} have been predominantly qualitative
or point-specific: the Earth Similarity Index \citep{Schulze-Makuch2011}
produces a single scalar per body with no spatial resolution;
\citet{Hendrix2019} reviewed icy ocean world habitability using qualitative
criteria matrices; \citet{Affholder2021} applied rigorous Bayesian inference
to Enceladus plume data but at a single spatial point and epoch.  None
integrates the full suite of Cassini data modalities into a spatially-resolved,
multi-epoch probabilistic framework.  A detailed methodological comparison is
given in Appendix~\ref{app:lit_review}.

This thesis develops a Bayesian Beta-conjugate posterior framework combining
eight observational proxy features from the complete Cassini dataset into a
habitability probability map at \SI{4490}{\metre\per\text{px}} resolution,
applied across five temporal epochs from the Late Heavy Bombardment ($-3.5\gya$)
to the red-giant future ($+5.9\gya$).  The Beta-conjugate approach is preferred
for its analytical tractability at 6.49 million grid pixels, its naturally
bounded $[0,1]$ output domain, and the explicit uncertainty quantification
provided by 95\% highest-density interval credible intervals at every pixel.

The specific research questions are:
\begin{description}
  \item[RQ1] \textit{What is the spatial distribution of present-day surface
    habitability, and which environments produce the highest scores?}

  \item[RQ2] \textit{How did habitability differ during the Late Heavy
    Bombardment, when impact-melt ponds provided transient liquid water?}

  \item[RQ3] \textit{How does habitability evolve across the temporal sequence
    from the LHB ($-3.5\gya$) to the red-giant ocean ($+5.9\gya$)?}

  \item[RQ4] \textit{Does the derived habitability map independently validate
    the \dragonfly{} landing site at Selk crater?}
\end{description}
